% =====================================================================
%  c06 —— PDF p.46–54（原书 p.40–48）
%  §3 Moduli stack 的相交理论；§3.1 由粘合构造叠
%  说明：交换图均用 tikz-cd 重画（原书图幅小、字密，裁图后放大即模糊）。
% =====================================================================

% ===== 原书 p.40 =====
\section{Moduli stack 的相交理论}

\subsection{由粘合构造叠}

为了用二次微分（quadratic differential）定义代数叠（algebraic stack）或与之相伴的解析叠（analytic stack）的胞腔分解，这里把叠实现为余等化子（coequalizer）。

设$C$为带 Grothendieck 拓扑$E$的范畴，$S$为$(C,E)$上的叠。$F:S\to C$为$S$的结构态射。$S$为$E$-叠，是指存在可表的叠$X_0$及$E$-满射的态射$\pi_0:X_0\to S$，且$S\to S\times_CS$可表。当$C$为$\C$上的代数簇、$E$为平展拓扑时，这无非是$\C$代数叠。由$\pi_0$的可表性，$X_1=X_0\times_SX_0$再次可表。令$pr_i:X_1=X_0\times_SX_0\to X_0$为自然投影，$\Delta:X_0\to X_1$为对角态射。

一般地，在范畴$C$中，两个对象$X_0,X_1$及三个态射$pr_1,pr_2:X_1\to X_0$、$\Delta:X_0\to X_1$的组合若满足$pr_1\circ\Delta=pr_2\circ\Delta=\id_{X_0}$，则称为$C$的 1-单纯对象（1-simplicial object），记作$X_1\rightrightarrows X_0$。进一步，1-单纯对象$X_1\rightrightarrows X_0$称为对称的，是指存在对合$\theta:X_1\to X_1$满足$pr_1=pr_2\circ\theta$。给定 1-单纯对象$X_1\rightrightarrows X_0$及$X_0$上的对象$a_0:Y\to X_0$时，若写$X_1\times_{X_0}Y$，除非另有说明，都表示把$X_1$经由$pr_2$看作$X_0$上的对象；若写$Y\times_{X_0}X_1$，则表示把$X_1$经由$pr_1$看作$X_0$上的对象。

$X_0$上的对象$Y$称为带下降数据（descent data），是指给定同构$\psi:X_1\times_{X_0}Y\to Y\times_{X_0}X_1$，且下图可换。
\begin{center}
\begin{tikzcd}[column sep=3.6em, row sep=2.8em]
X_1\times_{X_0}Y \arrow[r, "\psi"'] \arrow[d, "pr_1"'] & Y\times_{X_0}X_1 \arrow[d, "pr_2"]\\
X_1 \arrow[r, "\simeq"] & X_1
\end{tikzcd}
\end{center}

% ===== 原书 p.41 =====
下降数据称为满足上闭链条件（cocycle condition），是指下图可换。
\begin{center}
\begin{tikzcd}[column sep=4em, row sep=3.4em]
X_1\times_{X_0}X_1\times_{X_0}Y \arrow[r, "pr_{23}^*(\psi)"] \arrow[d, equal] & X_1\times_{X_0}Y\times_{X_0}X_1 \arrow[d, "pr_{12}^*(\psi)"]\\
X_1\times_{X_0}X_1\times_{X_0}Y \arrow[r, "pr_{13}^*(\psi)"] & Y\times_{X_0}X_1\times_{X_0}X_1
\end{tikzcd}
\end{center}

又，若给定$C$的对象$Y,\tilde Y,X_0$及三个态射$a_\#:\tilde Y\to Y$，$a_0:\tilde Y\to X_0$，则称$\tilde Y$为$(Y,X_0)$-双对象（biobject）。$(Y,X_0)$-双对象的平凡化（trivialization），是指同构$a_1:\tilde Y\times_Y\tilde Y\cong\tilde Y\times_{X_0}X_1$，且下图可换。
\begin{center}
\begin{tikzcd}[column sep=2.6em, row sep=3em]
\tilde Y & \tilde Y\times_Y\tilde Y \arrow[l, "pr_1"'] \arrow[r, "a_0\circ pr_2"] \arrow[d, "a_1"] & X_0\\
\tilde Y \arrow[u, equal] & \tilde Y\times_{X_0}X_1 \arrow[l, "pr_1"'] \arrow[r, "pr_2\circ pr_2"] & X_0 \arrow[u, equal]
\end{tikzcd}
\end{center}
以及下图也可换。
\begin{center}
\begin{tikzcd}[column sep=3.6em, row sep=2.8em]
\tilde Y\times_Y\tilde Y \arrow[r, "\theta_{\tilde Y}"'] \arrow[d] & \tilde Y\times_Y\tilde Y \arrow[d]\\
X_1 \arrow[r, "\theta"'] & X_1
\end{tikzcd}
\end{center}

此时，态射的合成
\[
X_1\times_{X_0}\tilde Y\xrightarrow{\theta\times1}\tilde Y\times_{X_0}X_1\xleftarrow{a_1}\tilde Y\times_Y\tilde Y\xrightarrow{\theta_{\tilde Y}}\tilde Y\times_Y\tilde Y\xrightarrow{a_1}Y\times_{X_0}X_1
\]
给出下降数据。它满足上闭链条件时，称此平凡化满足上闭链条件。又，取上述的$Y$为$X_0$的情形，把$X_1$经由$pr_1$看作$X_0=Y$上的对象，再经由$pr_2$赋予$a_0$-结构，则$X_1$可看作$(X_0,X_0)$-双对象。此时的平凡化$\psi:X_1\times_{pr_2,X_0,pr_1}X_1\to X_1\times_{pr_1,X_0,pr_1}X_1$称为 1-单纯对象的平凡化。若给定$X_1\rightrightarrows X_0$的对称化（simmetrization），且此平凡化满足上闭链条件，则称为满足上闭链条件的对称 1-单纯对象。

\transnote{原文此处的西文拼作「simmetrization」，同一页稍后又拼作「symmetrization」，此处照原文印出。}

回到$(C,E)$-叠的平展满射态射$X_0\to S$及$X_1=X_0\times_SX_0$的情形看看。$X_1$上的对称化由$X_0$的分量交换$(x_1,x_2)\mapsto(x_2,x_1)$给出。平凡化$X_1\times_{2,X_0,1}X_1\to X_1\times_{1,X_0,1}X_1$若按$(x_1,x_2,x_2,x_3)\mapsto$
% ===== 原书 p.42 =====
$(x_1,x_2,x_1,x_3)$确定，则容易验证它满足上闭链条件。因此这个 1-单纯对象是对称的，且满足上闭链条件。

一般地，对满足上闭链条件的对称 1-单纯对象$X_1\rightrightarrows X_0$，定义群胚中的范畴（category in groupoid）$\Coeq(X_1\rightrightarrows X_0)$。

\textbf{Object}：对象是如下三元组
\[
a=(a_\#:\tilde Y\to Y,\ a_0:\tilde Y\to X_0,\ a_1:\tilde Y\times_Y\tilde Y\to\tilde Y\times_{X_0}X_1),
\]
且满足
\begin{enumerate}[label=(\arabic*)]
\item 由$(a_\#,a_0)$，$\tilde Y$成为$(Y,X_0)$-双对象。
\item $a_\#$为平展满射。
\item $a_1$是此双对象结构的平凡化，且满足上闭链条件。
\end{enumerate}

\textbf{Morphism}：对两个对象$a^{(i)}=(a_\#^{(i)}:\tilde Y^{(i)}\to Y^{(i)},a_0^{(i)}:\tilde Y^{(i)}\to X_0,a_1^{(i)}:\tilde Y^{(i)}\times_{Y^{(i)}}\tilde Y^{(i)}\to\tilde Y^{(i)}\times_{X_0}X_1)$，态射$f:a^{(1)}\to a^{(2)}$是指两个态射$f_\#:Y^{(1)}\to Y^{(2)}$，$f_0:\tilde Y^{(1)}\to\tilde Y^{(2)}$的组，满足以下条件。
\begin{enumerate}[label=(\arabic*)]
\item 下图可换。
\begin{center}
\begin{tikzcd}[column sep=3.4em, row sep=2.8em]
\tilde Y^{(1)} \arrow[r, "f_0"] \arrow[d, "a_\#^{(1)}"'] & \tilde Y^{(2)} \arrow[d, "a_\#^{(2)}"]\\
Y^{(1)} \arrow[r, "f_\#"] & Y^{(2)}
\end{tikzcd}
\end{center}
\item $f_0$成为$X_0$上的态射。
\item 下图可换。
\begin{center}
\begin{tikzcd}[column sep=4.6em, row sep=3em]
\tilde Y^{(1)}\times_{Y^{(1)}}\tilde Y^{(1)} \arrow[r, "a_1^{(1)}"] \arrow[d, "(f_0\times f_0)"'] & \tilde Y^{(1)}\times_{X_0}X_1 \arrow[d, "(f_0\times1_{X_0})"]\\
\tilde Y^{(2)}\times_{Y^{(2)}}\tilde Y^{(2)} \arrow[r, "a_1^{(2)}"] & \tilde Y^{(2)}\times_{X_0}X_1
\end{tikzcd}
\end{center}
\item 诱导的态射$(a_\#,f_0):\tilde Y^{(1)}\to Y^{(1)}\times_{Y^{(2)}}\tilde Y^{(2)}$为同构。
\end{enumerate}

% ===== 原书 p.43 =====
至此对象与态射已定义完毕。取从$\Coeq(X_1\rightrightarrows X_0)$到$C$的函子$F$，为把上述$a$对应到$Y$的函子。此时，由态射的定义，它自然构成群胚中的范畴。

\begin{proposition}
$S\to C$为$(C,E)$-叠，$X_0\to S$为平展满射态射。令$X_1=X_0\times_SX_0$，则$\Coeq(X_1\rightrightarrows X_0)$与$S$作为群胚中的范畴等价。
\end{proposition}

\begin{proof}
先构造函子$F:S\to\Coeq(X_1\rightrightarrows X_0)$。给定$S$的对象$a$及$\pi(a)=Y$，与给定函子$a:Y\to S$等价，因此把$Y$看作$S$上的对象。由于$X_0\to S$可表，范畴的纤维积$Y\times_SX_0$可由$\tilde Y$表出。再由$pr_2$定义$\tilde Y\to X_0$。于是$\tilde Y$成为$(Y,X_0)$-双对象。又，对这个$\tilde Y$，由
\[
\tilde Y\times_Y\tilde Y\simeq(Y\times_SX_0)\times_Y(Y\times_SX_0)\simeq Y\times_S(X_0\times_SX_0)\simeq\tilde Y\times_{X_0}X_1
\]
的合成给出平凡化$a_1$。进而，由$pr_1:\tilde Y\to Y$，$pr_2:\tilde Y\to X_0$分别定义$a_\#,a_0$，并由$(a_\#:\tilde Y\to Y,a_0:\tilde Y\to X_0,a_1:\tilde Y\times_Y\tilde Y\to\tilde Y\times_{X_0}X_1)$定义$F(a)$。

容易知道这个平凡化满足上闭链条件。此外，这个函子与纤维群胚的结构相容。也就是说，给定$a^{(1)}=(a_\#^{(1)}:\tilde Y^{(1)}\to Y^{(1)},a_0^{(1)},a_1^{(1)})$及态射$f_\#:Y^{(2)}\to Y^{(1)}$时，$f_\#^*a$的$F$像就是$F(a)$在$Y^{(1)}$上的基变换。

其次定义函子$G:\Coeq(X_1\rightrightarrows X_0)\to S$。现在设$\Coeq(X_1\rightrightarrows X_0)$的对象$a=(a_\#,a_0,a_1)$已给定。此时由$\pi\circ a_0:\tilde Y\to S$定义$S_{\tilde Y}$的对象$\alpha=\pi\circ a_0$。由$pr_i:\tilde Y\times_Y\tilde Y\to\tilde Y$拉回，

% ===== 原书 p.44 =====
同构$f_{12}:pr_1^*(\alpha)\simeq pr_2^*(\alpha)$由下图可知。
\begin{center}
\begin{tikzcd}[column sep=3.2em, row sep=3em]
\tilde Y\times_Y\tilde Y \arrow[r, "pr_2"] \arrow[d, "a_1"'] & \tilde Y \arrow[r, equal] & \tilde Y \arrow[d, "a_0"]\\
\tilde Y\times_{X_0}X_1 \arrow[r, "pr_2"] \arrow[d, "pr_1"'] & X_1 \arrow[r, "pr_2"] \arrow[d, "pr_1"'] & X_0 \arrow[d, "\pi"]\\
\tilde Y \arrow[r, "a_0"] & X_0 \arrow[r, "\pi"] & S
\end{tikzcd}
\end{center}

由于$a_1$满足上闭链条件，可知$f_{12}$在$S_{\tilde Y\times_Y\tilde Y\times_Y\tilde Y}$上满足上闭链条件$pr_{12}^*f_{12}\circ pr_{23}^*f_{12}=pr_{13}^*f_{12}$。由叠的有效下降条件，存在$S_Y$的对象$\beta$，使$\alpha=a_\#^*\beta$。对$\Coeq(X_1\rightrightarrows X_0)$的对象$a$，使这个$\beta$对应过去，就定义了函子$G$。

此时，对$\beta:Y\to S$，$G\circ F(\beta)=\beta$。这由下图及$a_\#^*(\beta)=\pi\circ a_0$可知。
\begin{center}
\begin{tikzcd}[column sep=3em, row sep=2.8em]
Y\times_SX_0 \arrow[r] \arrow[d, equal] & X_0 \arrow[r, "\pi"] \arrow[d, equal] & S\\
\tilde Y \arrow[r, "a_\#"] & Y \arrow[r, "\beta"] & S
\end{tikzcd}
\end{center}

$F\circ G=1$可如下证明。取$a=(a_\#,a_0,a_1)$为$\Coeq(X_1\rightrightarrows X_0)$的对象，置$G(a)=(\beta:Y\to S)$。由于$F$与基变换相容，用$a_\#:\tilde Y\to Y$作基变换，得$a_\#^*(F(\beta))=F(a_\#^*\beta)=F(\pi\circ a_0)$，故$F(\pi\circ a_0)$由下图得到。
\begin{center}
\begin{tikzcd}[column sep=3em, row sep=2.8em]
\tilde Y\times_{X_0}X_1 \arrow[r] \arrow[d] & X_1 \arrow[r] \arrow[d] & X_0 \arrow[d]\\
\tilde Y \arrow[r] & X_0 \arrow[r] & S
\end{tikzcd}
\end{center}

另一方面，$a_\#^*(a)$由可换图式
\begin{center}
\begin{tikzcd}[column sep=3em, row sep=2.8em]
\tilde Y\times_Y\tilde Y \arrow[r] \arrow[d] & \tilde Y \arrow[d]\\
\tilde Y \arrow[r] & Y
\end{tikzcd}
\end{center}
得到，故$a_1$给出同构$a_\#^*(a)\simeq a_\#(F\circ G(a))$。这个同构是$(a_\#,a_0):\tilde Y\to Y\times_SX_0$的$a_\#$拉回。

% ===== 原书 p.45 =====
（这个态射由$\beta\circ a_\#=\pi\circ a_0$得到。）因此$(a_\#,a_0)$成为同构。因此$F\circ G(a)=a$。
\end{proof}

接下来考虑给定 1-单纯对象时构造$(C,E)$-叠。为此说明$E$-粘合（$E$-glewing）的条件，并叙述 Isom 函子。设$\tilde Z\to Z$，$Z\to Y$为两个$E$-满射态射。现在设$\tilde Z$关于$Z\times_YZ\rightrightarrows Z$带有下降数据，且满足上闭链条件。把这个下降数据记作$\psi:\tilde Z\times_YZ\to Z\times\tilde Z$。$E$-粘合封闭，是指在上述情形下，总存在除同构外唯一的对象$\tilde Y$及可换图式
\begin{center}
\begin{tikzcd}[column sep=3.4em, row sep=2.8em]
\tilde Z \arrow[r] \arrow[d] & \tilde Y \arrow[d]\\
Z \arrow[r] & Y
\end{tikzcd}
\end{center}
且自然诱导的态射$\tilde Z\to Z\times_Y\tilde Y$为同构。

其次，对$\Coeq(X_1\rightrightarrows X_0)_Y$的对象$a^{(1)}=(a_\#^{(1)}:\tilde Y^{(1)}\to Y^{(1)},a_0^{(1)},a_1^{(1)})$、$a^{(2)}=(a_\#^{(2)}:\tilde Y^{(2)}\to Y^{(2)},a_0^{(2)},a_1^{(2)})$，定义$(C/Y)$上的函子$\Isom_Y(a^{(1)},a^{(2)})$。给定$Z\to Y$时，把$\Isom_Y(a^{(1)},a^{(2)})(Z)$定义为$(Z,X_0)$上的双对象同构$h_0:\tilde Y^{(1)}\times_YZ\to\tilde Y^{(2)}\times_YZ$的集合，其中要求图式$(*)$可换。
\begin{center}
$(*)$\quad
\begin{tikzcd}[column sep=5.4em, row sep=3.4em]
(\tilde Y^{(1)}\times_YZ)\times_Z(\tilde Y^{(1)}\times_YZ) \arrow[r, "{(h_0,h_0)}"] \arrow[d, "a_1^{(1)}"'] & (\tilde Y^{(2)}\times_YZ)\times_Z(\tilde Y^{(2)}\times_YZ) \arrow[d, "a_1^{(2)}"]\\
(\tilde Y^{(1)}\times_{X_0}X_1)\times_YZ \arrow[r, "{(h_0,1)}"'] & (\tilde Y^{(2)}\times_{X_0}X_1)\times_YZ
\end{tikzcd}
\end{center}

\begin{remark}
当$(C,E)$为$(An,etale)$、$(PC^\infty,loc.iso)$、$(Top,loc.iso)$时，$E$-粘合封闭。
\end{remark}

\begin{theorem}
设给定满足上闭链条件的对称 1-单纯对象$X_1\rightrightarrows X_0$，且$pr_i:X_1\to X_0$为$E$-满射。进一步，设
\begin{enumerate}[label=(\arabic*)]
\item $(C,E)$关于$E$-粘合封闭。
\item 对任意$\Coeq(X_1\rightrightarrows X_0)_Y$的对象$a^{(1)},a^{(2)}$，$\Isom_Y(a^{(1)},a^{(2)})$可表。
\end{enumerate}
则$\Coeq(X_1\rightrightarrows X_0)$成为$(C,E)$-叠。
\end{theorem}

% ===== 原书 p.46 =====
\begin{proof}
记$\Coeq(X_1\rightrightarrows X_0)$为$S$。它成为群胚中的范畴，这一点已经看过。

\textbf{有效下降条件}：确认有效下降条件。设$Z\to Y$为平展满射态射，$a$为$S_Z$的对象。进一步设$pr_i:Z\times_YZ\to Z$的拉回$pr_1^*a,pr_2^*a$作为$S_{Z\times_YZ}$的对象带有同构$(h_\#,h_0:\tilde Z\times_Z(Z\times_YZ)\to(Z\times_YZ)\times_Z\tilde Z)$，且满足上闭链条件。先考虑构造$\tilde Y\to Y$。对$\tilde Z$已给定关于$pr_i:Z\times_YZ\to Z$的下降数据，且它满足上闭链条件。由于$(C,E)$粘合封闭，存在$\tilde Y\to Y$使$\tilde Z=\tilde Y\times_YZ$。其次看$a_0:\tilde Z\to X_0$是否经$\tilde Y$分解。为此只需验证$a_0\circ pr_1=a_0\circ pr_2$，这由$\tilde Y$的定义与下图可知。
\begin{center}
\begin{tikzcd}[column sep=3.4em, row sep=2.8em]
\tilde Z\times_Y\tilde Z \arrow[r, "\simeq"] \arrow[d, "\simeq"'] & \tilde Z\times_YZ \arrow[d]\\
Z\times_Y\tilde Z \arrow[r] & X_0
\end{tikzcd}
\end{center}

进而，同构
\[
(\tilde Y\times_Y\tilde Y)\times_YZ\simeq\tilde Z\times_Z\tilde Z\xrightarrow{a_1}(\tilde Y\times_{X_0}X_1)\times_YZ\simeq\tilde Z\times_{X_0}X_1
\]
限制到$Z\times_YZ$上时与$pr_1,pr_2$一致。这由上述同构来自下降数据$h_0$而知。

\textbf{$\Hom$ 的层条件}：设$Z\to Y$为平展满射，$a^{(1)},a^{(2)}$为$S_Y$的对象。此时
\[
0\to\Hom_{S_Y}(a^{(1)},a^{(2)})\to\Hom_{S_Z}(\pi_0^*a^{(1)},\pi_0^*a^{(2)})\rightrightarrows\Hom_{S_{Z\times_YZ}}(\pi_1^*a^{(1)},\pi_1^*a^{(2)})
\]
正合。这里$\pi_0:Z\to Y$，$\pi_1:Z\times_YZ\to Y$为自然态射。关于$a^{(1)},a^{(2)}$沿用此前的记号。现在设态射$h:\pi_0^*a^{(1)}\to\pi_0^*a^{(2)}$由$(h_\#,h_0)$给出。这里$h_0$是态射$h_0:\tilde Y^{(1)}\times_YZ\to\tilde Y^{(2)}\times_YZ$。

% ===== 原书 p.47 =====
设$pr_1^*h=pr_2^*h$，则$pr_1^*h_0=pr_2^*h_0$，故存在$\bar h_0:\tilde Y^{(1)}\to\tilde Y^{(2)}$使$\pi_0^*\bar h_0=h_0$。进而对这个$\bar h_0$，考虑图式
\begin{center}
\begin{tikzcd}[column sep=4.4em, row sep=3em]
\tilde Y^{(1)}\times_Y\tilde Y^{(1)} \arrow[r, "\bar h_0\times\bar h_0"] \arrow[d, "a_1^{(1)}"'] & \tilde Y^{(2)}\times_Y\tilde Y^{(2)} \arrow[d, "a_1^{(2)}"]\\
\tilde Y^{(1)}\times_{X_0}X_1 \arrow[r, "\bar h_0\times1"] & \tilde Y^{(2)}\times_{X_0}X_1
\end{tikzcd}
\end{center}
用$\pi_0$作基变换后可换，故由$E$-满射性，上图也可换。这样$\Hom_{S_Y}(a^{(1)},a^{(2)})$就构造出来了。唯一性由$E$-满射性可知。

\textbf{$X_0\to S$的平展满射性}：要证$X_0\to S$的平展满射性，只需说明$Y\times_SX_0$可由$\tilde Y$表出。$\tilde Y\xrightarrow{a_\#}Y\to S$与$\tilde Y\xrightarrow{a_0}X_0\to S$分别作为$S$的对象由下面的图式给出。
\begin{center}
\begin{tikzcd}[column sep=3.2em, row sep=2.8em]
\tilde Y\times_Y\tilde Y \arrow[r] \arrow[d] & \tilde Y \arrow[r] \arrow[d] & X_0\\
\tilde Y \arrow[r] & Y
\end{tikzcd}
\end{center}
\begin{center}
\begin{tikzcd}[column sep=3.2em, row sep=2.8em]
\tilde Y\times_{X_0}X_1 \arrow[r] \arrow[d] & X_1 \arrow[r] \arrow[d] & X_0\\
\tilde Y \arrow[r] & X_0
\end{tikzcd}
\end{center}
诱导的$\tilde Y\times_Y\tilde Y\to\tilde Y$与$\tilde Y\times_{X_0}X_1\to\tilde Y$经由$a_1$成为同构。由此得到态射$\tilde Y\to Y\times_SX_0$。反之，给定$Z\to Y$，$Z\to X_0$及同构$h_0:Z\times_Y\tilde Y\to Z\times_{X_0}X_1$时，利用$pr_2:Z\times_Y\tilde Y\to\tilde Y$，$pr_2:Z\times_{X_0}X_1\to X_1$，得到态射$Z\times_Y\tilde Y\to\tilde Y\times_{X_0}X_1$。合成$\tilde Y\times_{X_0}X_1\simeq\tilde Y\times_Y\tilde Y$所得的态射$Z\times_Y\tilde Y\to\tilde Y\times_Y\tilde Y$可下降到$Z\to X_1$，这与前面同样可证。可知它给出上述对应的逆对应，从而$Y\times_SX_0\simeq\tilde Y$。

\textbf{对角态射的可表性}：现在设$a^{(1)},a^{(2)}$作为两个$S_Y$的对象已给定。也就是说，给定$Y\to S\times_CS$。此时范畴的纤维积$S\times_{S\times_CS}Y$是$Z$上的对象与$(Z,X_0)$-双对象同构的组$(Z\to Y,h_0:\tilde Y^{(1)}\times_YZ\simeq\tilde Y^{(2)}\times_YZ)$且图式$(*)$可换者组成的范畴。由假设，$\Isom_Y(a^{(1)},a^{(2)})$可表，故它与$\Isom_Y(a^{(1)},a^{(2)})$上$C$的对象组成的范畴等同。由此可知对角态射可表。
\end{proof}

% ===== 原书 p.48 =====
其次，设给定满足定理 3.2 条件的 1-单纯对象$X_1\rightrightarrows X_0$、$X'_1\rightrightarrows X'_1$及它们之间的态射$(X'_1\rightrightarrows X'_0)\to(X_1\rightrightarrows X_0)$，定义映射$\Coeq(X'_1\rightrightarrows X'_0)\to\Coeq(X'_1\rightrightarrows X'_0)$，并叙述它为同构的条件。

\transnote{原书此处印作$X'_1\rightrightarrows X'_1$，且所定义映射的两端都印作$\Coeq(X'_1\rightrightarrows X'_0)$；由文脉当为$X'_1\rightrightarrows X'_0$及$\Coeq(X'_1\rightrightarrows X'_0)\to\Coeq(X_1\rightrightarrows X_0)$。今照原文翻译。}

\begin{theorem}
设给定满足定理 3.2 条件的 1-单纯对象$X_1\rightrightarrows X_0$、$X'_1\rightrightarrows X'_1$及它们之间的态射$\phi:(X'_1\rightrightarrows X'_0)\to(X_1\rightrightarrows X_0)$。此时$\Coeq(\phi):\Coeq(X'_1\rightrightarrows X'_0)\to\Coeq(X'_1\rightrightarrows X'_0)$自然定义。进一步，若由 1-单纯对象的映射自然诱导的映射
\[
X'_1\to X'_0\times_{X_0}X_1\times_{X_0}X'_0
\]
为同构，则$\Coeq(\phi)$为同构。
\end{theorem}

\begin{proof}
先对$\Coeq(X'_1\rightrightarrows X'_0)$的对象$\tilde Y'\to Y$，定义与之对应的$\Coeq(X_1\rightrightarrows X_0)$的函子$F$。为此，在$\tilde Y'$上的对象$\tilde Y'\times_{X_0}X_1$上，给定关于满足上闭链条件的$\tilde Y'\times_Y\tilde Y'\rightrightarrows\tilde Y'$的下降数据，构造$\tilde Y\to Y$。这里把$\tilde Y'\to X'_0\to X_0$的合成看作$\tilde Y'$为$X_0$上的对象。

只需定义同构
\[
\psi:(\tilde Y'\times_Y\tilde Y')\times_{pr_1,X_0}X_1\to(\tilde Y'\times_Y\tilde Y')\times_{pr_2,X_0}X_1
\]
即可。由态射$\tilde Y'\times_Y\tilde Y'\to\tilde Y\times_{X_0}X_1\to X_1$，把$\tilde Y'\times_Y\tilde Y'$看作$X_1$上的代数簇。此时
\begin{align*}
(\tilde Y'\times_Y\tilde Y')\times_{pr_1,X_0}X_1&\simeq(\tilde Y'\times_Y\tilde Y')\times_{X_1,pr_1}(X_1\times_{pr_1,X_0}X_1),\\
(\tilde Y'\times_Y\tilde Y')\times_{pr_2,X_0}X_1&\simeq(\tilde Y'\times_Y\tilde Y')\times_{X_1,pr_1}(X_1\times_{pr_2,X_0}X_1),
\end{align*}
故只需给出作为右$X_1$对象的同构$X_1\times_{pr_1,X_0,pr_1}X_1\to X_1\times_{pr_2,X_0,pr_1}X_1$。现在 1-单纯对象$X_1\rightrightarrows X_0$带有下降数据$X_1\times_{pr_1,X_0,pr_1}X_1\to X_1\times_{pr_2,X_0,pr_1}X_1$，故用它。由于$C$关于$E$-粘合封闭，$\tilde Y\to Y$构造出来了。

% ===== 原书 p.49 =====
进一步，在后半条件下，从$\Coeq(X_1\rightrightarrows X_0)$到$\Coeq(X'_1\rightrightarrows X'_0)$的函子，对$\Coeq(X_1\rightrightarrows X_0)$的对象$\tilde Y\to Y$，用$\tilde Y'=\tilde Y\times_{X_0}X'_0\to\tilde Y\to Y$给出，则
\begin{align*}
\tilde Y'\times_Y\tilde Y'&\simeq X'_0\times_{X_0}\tilde Y\times_Y\tilde Y\times_{X_0}X'_0\\
&\simeq X'_0\times_{X_0}\tilde Y\times_{X_0}X_1\times_{X_0}X'_0\\
&\simeq\tilde Y\times_{X_0}X'_0\times_{X_0}X_1\times_{X_0}X'_0\\
&\simeq(\tilde Y\times_{X_0}X'_0)\times_{X'_0}X'_1
\end{align*}
给出$\Coeq(X'_1\rightrightarrows X'_0)$的对象。这些函子给出$\Coeq(X_1\rightrightarrows X_0)$与$\Coeq(X'_1\rightrightarrows X'_0)$的范畴等价。
\end{proof}

由定理 3.1、定理 3.2，对$(Sch/\C,etale)$-叠$S$，如下定义$(An,etale)$-叠$S^{an}$。先取一个平展满射态射$X_0\to S$，定义$X_1=X_0\times_SX_0$，再考虑$X_1^{an}\rightrightarrows X_0^{an}$，则$X_1^{an}\rightrightarrows X_0^{an}$成为满足上闭链条件的对称 1-单纯对象。进一步，由于$pr_i:X_1^{an}\to X_0^{an}$为平展满射，对$Y\in(An)$及$a^{(1)},a^{(2)}\in\Coeq(X_1^{an}\rightrightarrows X_0^{an})_Y$，可知$\Isom_Y(a^{(1)},a^{(2)})$可表。由此$S^{an}=\Coeq(X_1^{an}\rightrightarrows X_0^{an})$成为$(An,etale)$-叠。又，由定理 3.3 可知$S^{an}$与平展满射态射$X_0\to S$的取法无关。称$S^{an}$为$S$的相伴解析叠（associated analytic stack）。完全同样地，对$(An,etale)$-叠，也可对应$(PC^\infty,loc.iso)$-叠及$(Top,loc.iso)$-叠。
